\scanpage{48}
\mspage{1}
% archive foliation: 50
\unsure{10369CNU|10369GNU}\qquad 15.1.

$X$ \deleted{non singulière} \inserted{lisse de type $g,\nu$} (sur $\overline{\mathbb Q}$)

$Y$ stable non lisse type $(g,\nu)$ avec propagateur normal $N$ d'où un groupe fondamental \unsure{ét.}\deleted{\unsure{indép.}} $\unsure{\Pi|\pi}_1([Y,N])$

\diagram{\begin{array}{c}X\overset f{\underset T\simeq}Y\end{array}} un isomorphisme de Teichmüller.

\uline{NB} Cet isomorphisme induit un isom de groupes extérieurs \deleted{\unsure{celui est de}} (\unsure{il est déterminé} plus bas)
\diagram{\begin{tikzcd}[column sep=large]\unsure{\Pi|\pi}_1(X)\arrow[r]&\unsure{\Pi|\pi}_1([Y,N])\end{tikzcd}}

Quand on se donne un ``pt base sur $X$'' (ordinaire ou à l'infini), et un sur $[Y,N]$ (alors --- et de plus distinct des pts singuliers de $Y$) on a une notion d'isom de Teichmüller respectant les pts marqués, \deleted{et \unsure{d'un}} (on suppose qu'ils sont tous deux à distance finie, ou tous deux à l'infini). Et on trouve alors un isom de groupes (pas ext. !)
\diagram{\begin{tikzcd}[column sep=large]\unsure{\Pi|\pi}_1(X,\unsure{\bullet|x})\arrow[r]&\unsure{\Pi|\pi}_1([Y,N],y)\end{tikzcd}}

En tout état de cause, \deleted{\unsure{les}} l'isom extérieur déjà n'est pas compatible avec les actions ext. de $\operatorname{Gal}(\overline{\mathbb Q}/K)$, $K$ \inserted{\unsure{est}} une extension finie de $\overline{\mathbb Q}/\mathbb Q$ assez grande

\margin{\uline{NB} Mettre en évidence \uline{trois} classes de conjugaison d'\deleted{éléments} hom. $\widehat{\mathbb Z}(1)\longrightarrow\unsure{\Pi|\pi}_1(X)$ (deux seulement si $\nu=0$, l'une des trois est alors ``le vide'').}
